% Source: 13_Chapter_6_The_Rule_of_Books_The_Formation_of_the_B.txt
\bookchapter[The Rule of Books]{CHAPTER 6}{THE RULE OF BOOKS}{The Formation of the Bible}

\leadin{D}{URING THE LAST GREAT} persecution of Christians in the Roman Empire, early in the fourth
century, a believer in Sicily was brought before the governor. He was charged with possessing a copy
of the Gospels.

“Where did these come from?” asked the judge, pointing to the books. “Did you bring them from your
home?”

“I have no home,” replied the prisoner, “as my Lord Jesus knows.”

Once again pointing to the Gospels, the judge said, “Read them!”

The Christian opened the Gospels and read, “Blessed are those who are persecuted for righteousness’
sake: for theirs is the kingdom of heaven.”

He turned to another place and read again, “If any man will come after me, let him deny himself, and
take up his cross, and follow me.”

That was too much. The judge ordered his prisoner away—to death.

Roman officials came to see that the suppression of Christianity demanded the destruction of the
Scriptures. So the last great persecution of Christians included the burning of the Scriptures.

To this day we find it almost impossible to think of the Christian faith without the Bible. It is
the foundation of Christianity’s evangelism, its teaching, its worship, and its morality. When we
look back over Christian history, we find few if any decisions more basic than those made during the
first three centuries surrounding the formation of the Bible. The Scriptures served not only as the
inspiration for believers facing martyrdom, but as the supreme standard for the churches threatened
by heresy. If catholic Christianity was orthodox, the Bible made it so, for the constant test of any
teaching was, what do the Scriptures say?

We need to ask, then, how did we get the Bible?

\section{THE BASICS OF THE BIBLE}

The name itself—Bible—suggests that Christians consider these writings special. Jerome, the
fourth-century translator, called them “the Divine Library.” He wanted to stress that the many books
were, in fact, one. Greek-speaking believers made the same point when they shifted from the early
plural form Biblia, meaning The Books, to The Bible, meaning The Book. Long before, Jews had faced
the same problem when they spoke of The Scriptures and Scripture. That explains how, in time, the
Bible and Scripture came to mean the same thing in Christian circles, the sixty-six books that
Christians consider the written Word of God.

Today, we find the Scriptures grouped under Old Testament (or Covenant) and New Testament. In the
ancient world a testament, or more often a covenant, was the term for a special relationship between
two parties. Occasionally we still speak of the marriage covenant, which binds husband and wife to
each other.

Used in the Bible, the term stands for the special relationship between God and man, initiated and
sustained by the grace of the Lord God. The old covenant was first between the Lord and Abraham,
then between God and Abraham’s descendants, the children of Israel. Later years knew them as Jews.
So the Old Testament contains the books that tell the story of the Jews and their ancient worship of
God.

Early Christians believed that Jesus of Nazareth was God’s promised Messiah, who established a new
covenant with his new people, the church. So the New Testament stands for the books telling the
story of Jesus Christ and the birth of the church.

The Bible contains two portions: the Old Testament, which the early Christians claimed along with
the Jews, and the New Testament, which the early Christians produced in spite of the Jews. The Old
Testament promised; the New Testament fulfilled.

The word for the special place these books occupy in Christianity is canon. The term from the Greek
language originally meant a measuring rod or, as we might say, a ruler. Eventually the word was
applied to the Scriptures: the books or listing of books that the church discerned were
authoritative or inspired by God. Since the first Christians were all Jews, Christianity was never
without a canon, or, as we say, Scripture. Jesus himself clearly accepted the Old Testament as God’s
word to man. “Scripture cannot be broken,” he said. “Everything written about me in the law of Moses
and the prophets and the psalms must be fulfilled” (John 10:35; Luke 24:44).

Jesus believed the statements of Scripture, endorsed its teaching, obeyed its commands, and set
himself to fulfill the pattern of redemption it laid down. Early Christians were simply heirs of
this attitude. Had not the hopes and plans of the old covenant come true in Jesus? Had not the
promised messianic age dawned in him?

Christians, almost without exception, embraced the Old Testament as their own. Early Christians,
both Jews and Gentiles, believed that Jesus was the fulfillment of the promises running through the
Old Testament. Christians also inherited strategies for interpreting from Jesus. Jews in Alexandria
had read much of the Bible as allegory. They sought a deeper spiritual or intellectual message
beyond the literal meaning. Still others employed a method called typology. The reader discerns a
pattern (type) or correspondence between two images or stories. Thus, the story of Jesus may be read
with an awareness that he is like an Old Testament sacrifice or an Old Testament deliverer such as
Moses.

By the third century the church had sophisticated scholars who could defend the Christian claim to
the Old Testament by the use of allegory. The most influential was a teacher at Alexandria named
Origen, who spoke of the different levels of Scripture:

\begin{blockquote}
The Scriptures were composed through the Spirit of God, and have both a meaning which is obvious,
and another which is hidden from most readers~.~.~.~. The whole law is spiritual, but the inspired
meaning is not recognized by all—only by those who are gifted with the grace of the Holy Spirit in
the word of wisdom and knowledge.
\end{blockquote}

Christian appeals to allegory infuriated pagan critics of the faith because their case depended on
their taking the Old Testament at face value. The move remained popular, however, since it enabled
Origen and other believers to find the Christian message just beneath the surface of the Old
Testament.

\section{THE QUESTION OF THE APOCRYPHA}

When Christians retained the Old Testament for their own use, they did not settle completely just
which books this included. To this day Christians differ over the inclusion or rejection of the
so-called Apocrypha in the Old Testament list of books. The term stands for twelve or fifteen books,
depending upon how you group them, that Roman Catholics and Eastern Orthodox accept as canonical and
most Protestants reject.

Some Jews embrace a canon that includes the typical Protestant canon for the Old Testament of
thirty-nine books. Scholars often picture this being more typical for the Jews in Palestine. They
numbered these books in clusters and organized them by a three-fold division of Law, Prophets, and
Writings. The three-fold structure of this emerging canon may explain Jesus referring to the Law of
Moses, the prophets, and the Psalms (the first book in the Writings, Luke 24:44). Most scholars
believe Jesus quoted only from this smaller canon, which is frequently but perhaps deceptively
called the “Palestinian Canon.”

The Greek translations of the Old Testament included the books of the Hebrew canon, with additional
books as well. Protestants routinely call these additional books the Apocrypha, or more generically
the Deuterocanonical books. The Greek translations, the chief of these being the Septuagint, were
influential in making the books of the Apocrypha more popularly circulated. Greek translations of
the Old Testament were made in the City of Alexandria. This larger canon, consisting of the Hebrew
canon and the Apocrypha, is called the “Alexandrian Canon.”

Early Christians differed over the question of the Apocrypha. In the West the influential Augustine,
the well-known bishop of Hippo, embraced the Apocrypha as part of the canon of Scripture. During the
sixteenth-century Reformation most Protestants rejected the Apocrypha as canonical. The Roman
Catholic Church, following Augustine, accepted the books. And that is how the churches differ to
this day.

From the beginning, however, Christians had more than the Old Covenant as their rule for faith.
During Jesus’ life on earth they had the Word made flesh, and after Jesus’ departure they had the
living leadership of the apostles. The reverence for the apostles’ message, whether oral or written,
as the authentic channel to the will of the Lord Jesus, is reflected throughout early Christian
literature.

During the days of the apostles, congregations often read letters from the companions of the Lord.
Some of these letters were obviously intended to be read in public worship, probably alongside some
portion of the Old Testament or with some sermon.

Churches also relied on accounts about the life of the Lord Jesus. The first gospels were not
written before AD 60 or 70, but their contents were partly available in written form before this.
Luke tells us that many had undertaken some account of the events of the life of Jesus.

The question is, out of this growing body of Christian literature, how did the twenty-seven books we
know as the New Testament come to be set apart as Scripture? How and when did they cross the line
between books regarded as important and even authoritative, and books regarded as holy and the Word
of God? To put it in one word, how did they become canonical?

Several factors were at work in this process. Some were internal characteristics of the developing
life of the churches; others were external, threats to the gospel arising from historical events and
pagan influences.

First, the books that are Scripture and are truly the Word of God have about them a self-evidencing
quality. They carry their uniqueness on their face. They have always exercised, and still exercise,
a transforming power upon the lives of men.

For example, as a young man Justin Martyr searched energetically for truth in a variety of
philosophical schools: first as a Stoic, then a Pythagorean, then a Platonist. But none of them
satisfied him. One day, while meditating alone by the seashore, perhaps at Ephesus, he met an old
man. During their conversation the stranger exposed the weaknesses of Justin’s thinking and urged
him to turn to the Jewish prophets. By reading Scripture, Justin became a Christian. Scores of other
men and women in the early days of the church had a similar experience: Tatian, Theophilus, Hilary,
Victorinus, Augustine.

One of the primary reasons, then, behind the adoption of the New Testament books as Holy Scripture
was this self-authenticating quality.

Second, certain Christian books were added to Scripture because they were used in Christian worship.
Even in the New Testament itself there are signs that the reading of Scripture was very much a part
of Christian congregational life. The apostle Paul urged the Colossians, “After this letter has been
read to you, see that it is also read in the church of the Laodiceans and that you in turn read the
letter from Laodicea” (Col. 4:16, NIV).

Justin Martyr, writing in the middle of the second century, gives us the first description of a
Christian service:

\begin{blockquote}
On the day called the Day of the Sun all who live in cities or in the country gather together to one
place, and the memoirs of the apostles or the writings of the prophets are read, as long as time
permits; then, when the reader has ceased, the president verbally instructs, and exhorts to the
imitation of these good things. Then we all rise together and pray.
\end{blockquote}

\noindent Thus we see by Justin’s time The Memoirs of the Apostles, which was his title for the
Gospels, were a central part of Christian worship.

The mere act of reading a book in Christian worship did not assure the writing an eventual place in
the canon. We know, for example, that Clement, Bishop of Rome, wrote a letter to the church at
Corinth about AD 96 and eighty years later it was still the custom in Corinth to read Clement’s
letter at public worship. Yet Clement’s letter was never added to the canon. Books read at the
worship of the church had a special position and had started on the road that led to entrance into
the canon of Scripture—but some did not make it.

Third, and perhaps the fundamental factor behind a Christian book’s acceptance into the New
Testament, was its ties to an apostle. This was the test of a book’s validity: was it written by an
apostle, or at least by a man who had direct contact with the circle of the apostles?

In the early church the apostles held a place that other men simply could not fill. Early believers
always regarded them as men who had a unique relationship with the Lord. Did not Jesus say, “He who
receives you receives me” (Matt. 10:40)?

Clement of Rome reflects this general attitude of Christians when he writes, “The apostles were made
evangelists to us by the Lord Christ; Jesus Christ was sent by God. Thus Christ is from God, and the
apostles from Christ~.~.~.~. The Church is built on them as a foundation” (1 Clement 42). Any gospel
or letter, therefore, that could make a strong claim to apostolic authorship stood a good chance of
acceptance as Scripture.

\section{A LIST OF CHRISTIAN BOOKS}

Given enough time the churches, under the influence of these factors, and perhaps others, probably
would have drawn up a list of canonical Christian writings. But certain events forced the hand of
the churches.

About AD 140 a wealthy and much-traveled ship owner from Sinope on the Black Sea came to Rome. His
name was Marcion. Although the son of a bishop, Marcion fell under the spell of the gnostic teacher
Cerdo, who believed that the God of the Old Testament was different from the God and Father of the
Lord Jesus Christ. The God of the Old Testament, he said, was unknowable; the Christian God had been
revealed. The Old Testament God was sheer justice; whereas the God of the New Covenant was loving
and gracious.

Marcion developed Cerdo’s distinction. He held that the Old Testament God was full of wrath and the
author of evil. This God, he said, was only concerned for the Jewish people. He was prepared to
destroy all other people. In contrast, the Christian’s God was a God of grace and love for all, who
disclosed himself in Jesus Christ, his Son.

Because he believed that the God of the Old Testament loved the Jews exclusively, Marcion rejected
the entire Old Testament and also those New Covenant writings that he thought favored Jewish
readers—for example Matthew, Mark, Acts, and Hebrews. He also rejected other Christian writings that
appeared to him to compromise his own views, including the Pastoral Letters (1 and 2 Timothy and
Titus). So he was left with only a mutilated version of Luke’s gospel (we suppose he omitted the
nativity stories) and ten letters of Paul. The Apostle to the Gentiles, it seems, was the only
apostle who did not corrupt the gospel of Jesus.

Marcion’s garbled Christian views were firmly repudiated by the church in Rome, and Marcion was
excommunicated from the church in AD 144. Before long, however, Marcionite churches appeared,
modeled on orthodox congregations. They had their own ministers and rituals. For example, they did
not use wine at communion, as a result of the ascetic emphasis of their teaching. Some of the
Marcionite beliefs spilled over into the various gnostic sects, and Marcionites were themselves
affected by gnostic views. Their ideas spread, however, throughout Italy and as far afield as
Arabia, Armenia, and Egypt. In the East they exercised a considerable influence for many decades. A
number of Marcionite villages existed near Damascus as late as the fourth century.

Most importantly, however, Marcion presented the orthodox churches with a twofold problem: his list
of New Testament books, shaped in the image of Paul, and his rejection of the Old Testament.

Marcion’s worship of Paul was little short of idolatry. As he saw it, Paul was the great enemy of
the law and the great spokesman for the gospel. He was in fact the supreme figure in the church.
Marcion believed Christ had descended from heaven twice, once to suffer and to die, and once to call
Paul and to reveal to Paul the true significance of his death. In heaven, said Marcion, Paul sits at
the right hand of Christ, who sits at the right hand of God.

As the North African lawyer Tertullian put it, Paul had become the apostle embraced by the heretics!
Of course, Marcion had to misinterpret Paul to make the apostle fit his beliefs, but that didn’t
make the churches’ problem any less real: how could they accept Paul’s letters as God’s word without
endorsing Marcionite teaching?

In the end Paul meant too much to the church to dismiss him because of Marcion’s extreme views. The
apostle’s letters were too well known and too widely used to discard them. The church chose,
instead, to restore the Pastorals and the letters of the other apostles and to link all the letters
to four gospels by using the book of Acts as the bridge. While the church treasured the grace of God
preached by Paul, it realized that jettisoning the Old Testament was suicidal. Does the New Covenant
make sense without the Old?

By retaining the Old Testament the church scored two important points. First, it insisted that faith
for the Christian would identify the idea of the creating God with the idea of the redeeming God.
Marcion’s message was too easy. Marcion had not only misread the Old Testament; he also broke the
unity seen in the entire Christian Scriptures: the same God who made the world also chose Israel;
that God sought to reclaim his creation by Jesus, who fulfills Israel’s destiny. Second, by
retaining the Old Testament the church underscored the importance of history for the Christian
faith. Christianity is a historical religion, not just in the sense that it comes from the past or
that it is associated with a historical character named Jesus. It is historical because it stems
from the belief that within history itself, in a particular place, at a particular time, God himself
took a hand in human affairs. And that means that living by faith for the Christian includes facing
the puzzles of human existence—all of the “why, Lord?”s of life—and still believing that God has
some good in mind.

If Marcion, a heretic, nudged the churches into thinking about forming a New Testament, another
troublemaker, Montanus, forced the churches into thinking about closing it.

\section{FRESH VOICES FROM GOD}

Christianity has always been a religion of the Spirit. According to the Fourth Gospel Jesus had
promised to his people the Paraclete, the Spirit of Truth, to guide them (John 16:13–15). How, then,
did there ever come a time when the church declared that all the inspired books that could be
written had been written, and that nothing more could ever be added to the written word of God? How
did it come about that, as Tertullian bitterly put it, “the Holy Spirit was chased into a book”?

In the second half of the second century a change was coming over the church. The days of enthusiasm
were passing and the days of ecclesiasticism were arriving. The church was no longer a place where
the Spirit of prophecy could be heard. More and more people were joining the churches, but the
distinction between church and world was fading. The church was becoming secularized; it was coming
to terms with heathen thought and culture and philosophy. The way of the cross was no longer rough
and steep.

Into this situation, sometime between AD 156 and 172, Montanus appeared, a voice in the wilderness
of Asia Minor. He came with a demand for a higher standard and a greater discipline and sharper
separation of the church from the world. Had he halted there, he could have done little but good,
but he went much further. He and his two prophetesses, Prisca and Maximilla, went about prophesying
in the name of the Spirit, and foretelling the speedy second coming of Christ. That in itself was
not extraordinary. But these new prophets, in contrast to prophets in biblical times, spoke in a
state of ecstasy, as though their personalities were suspended while the Paraclete spoke in them.
Montanus was convinced that he and his prophetesses were the God-given instruments of revelation,
the lyres across which the Spirit swept to play a new song. With that Montanus’s super spirituality
went too far.

Clearly the church had to act. One problem was simply disorder. Montanus as a herald of a new
spiritual vitality and a new challenge to holiness was one thing; but when Montanists insisted that
opposition to the new prophecy was blasphemy against the Holy Spirit, many churches split over the
question.

Montanus argued that the new age of the Spirit had displaced the previous ages. The Old Testament
age, with its ten commands or law, had been surpassed by the age of the Son with its more demanding
Sermon on the Mount. The age of the Spirit would call for even greater rigorous obedience, as well
as surpassing revelations. In the name of the Spirit, Montanus denied that God’s decisive and
normative revelation had occurred in Jesus Christ.

In the face of this challenge how could the church keep the gospel central? It had to make all later
Christian worship, teaching, and life center in Christ and the apostolic witness. Free utterances of
the Spirit would not guarantee that; Montanism was making this clear. The best way to make the
original apostolic gospel basic was to set apart the apostolic writings as uniquely authoritative.
This would require all later faith and action to be judged in the light of that central message.

It was not that the church had ceased to believe in the power of the Holy Spirit. The difference was
that in the first days the Holy Spirit had enabled men to write the sacred books of the Christian
faith; in the later days the Holy Spirit enabled men to understand, to interpret, and to apply what
had been written.

One of the reasons we know that the church assumed this position lies in the appearance of lists of
New Testament books. One of the first is a document written about AD 190. We call it the Muratorian
Canon, from its discoverer L. A. Muratori, who first published it in 1740. The document is damaged
at the beginning, and actually begins with Luke, but its list of books is as follows: Matthew, Mark,
Luke, John, Acts, 1 and 2 Corinthians, Ephesians, Philippians, Colossians, Galatians, 1 and 2
Thessalonians, Romans, Philemon, Titus, 1 and 2 Timothy, Jude, 1 and 2 John, the Apocalypse of John
(that is, the Revelation), the Apocalypse of Peter, and The Wisdom of Solomon. The last two, we
know, did not remain on the approved list. But by AD 190 the churches clearly accepted the idea of
Christian Scriptures alongside Jewish Scriptures, one fulfilling what the other promises.

By the early third century only a handful of books continued to create any question. Hebrews faced
some opposition in the western regions of the empire and Revelation was unpopular in the east. At
the outset of the fourth century, Eusebius, the church historian, summed up the situation and
indicated that James, 2 Peter, 2 and 3 John, and Jude were the only books spoken against by some
while recognized by others. Revelation, however, continued to bewilder him.

\figtitle{The New Testament as It Gained Acceptance by the Early Church}

The first complete list of books, as we have them today, came in an Easter letter written in 367 by
Bishop Athanasius from Alexandria. Shortly thereafter councils in North Africa at Hippo (393) and at
Carthage (397) published the same list.

\begin{sidebar}{The Logic of the Canon}
Everett Ferguson has expounded the process and progress of canonical thinking. First, believers come
to recognize that there are books that carry divine authority, and once the concept of “sacred
writings” is embraced, the process of canonical thinking is underway. Second, one begins to think of
a boundary for the canon. Sacred writings are different from other writings, and one might even
think there is a limited number. Third, believers might move to a closed canon, an explicit listing
of all the sacred writings. Finally, a consensus emerges around this common listing of sacred books.
\end{sidebar}

In one sense, of course, Christians created the canon. Their decisions concerning the books were a
part of history. In another sense, however, they were only recognizing those writings that had made
their authority felt in the churches. The shape of the New Testament shows that the early churches’
primary aim was to submit fully to the teachings of the apostles. In that purpose they shaped the
character of Christianity for all time. The faith remained catholic precisely because it was
apostolic.

\readinghead{Suggestions for Further Reading}
\begin{readinglist}
\entry Blackman, E. C. Marcion and His Influence. London: S.P.C.K., 1948.
\entry Bruce, F. F. The Canon of Scripture. Downers Grove, IL: InterVarsity, 1988.
\entry Campenhausen, Hans von. The Formation of the Christian Bible. Philadelphia: Fortress Press,
  1968.
\entry Filson, Floyd Vivian. Which Books Belong in the Bible? A Study of the Canon. Philadelphia:
  Westminster, 1957.
\entry Grant, Robert M. The Bible in the Church. New York: Macmillan, 1948.
\entry Metzger, B. M. The Canon of the New Testament. Oxford: Oxford University Press, 1987.
\entry *McDonald, Lee Martin. Formation of the Bible: The Story of the Church’s Canon. Peabody, MA:
  Hendrickson, 2012.
\end{readinglist}
